\section*{KẾT LUẬN}
\addcontentsline{toc}{section}{KẾT LUẬN}

Thông qua quá trình nghiên cứu và vận dụng lý thuyết môn học \textbf{``Khoa học Quản lý và Quản trị''} vào trường hợp điển hình của \textbf{Công ty Cổ phần FPT}, nhóm nghiên cứu rút ra những kết luận quan trọng sau:

Thứ nhất, thành công vượt trội của FPT trên thị trường trong nước và quốc tế suốt gần 4 thập kỷ qua bắt nguồn từ sự phân định rõ ràng và phối hợp hài hòa giữa hai trụ cột: \textbf{Quản trị công ty (Corporate Governance)} và \textbf{Quản lý điều hành (Management)}. Việc chủ động tách rời chức danh Chủ tịch HĐQT và Tổng Giám đốc từ năm 2013, kết hợp duy trì cơ cấu thành viên HĐQT độc lập chất lượng cao là minh chứng cho tư duy quản trị hiện đại, tiệm cận các chuẩn mực OECD và thông lệ quản trị quốc tế tốt nhất.

Thứ hai, mô hình quản lý của FPT là sự kết hợp sáng tạo giữa \textbf{Cơ cấu Công ty Mẹ -- Con (Holding)} và \textbf{Cơ cấu Ma trận đa chiều theo Đơn vị kinh doanh chiến lược (SBU) và Địa bàn hoạt động}. Mô hình này vừa đảm bảo sự kiểm soát tập trung về tầm nhìn chiến lược, tài chính và văn hóa tại công ty mẹ, vừa trao quyền tự chủ linh hoạt tối đa cho các công ty thành viên phản ứng nhanh với thị trường toàn cầu.

Thứ ba, cuộc cách mạng Trí tuệ Nhân tạo (AI) đang tạo ra bước ngoặt lịch sử cho FPT. Doanh nghiệp không chỉ cung cấp các giải pháp AI ra thị trường quốc tế (thông qua FPT.AI, akaBot, AI Factory với NVIDIA) mà còn đưa AI vào sâu trong từng mắt xích quản trị và quản lý nội bộ: từ tự động hóa quy trình, tăng 20--30\% năng suất lập trình, tối ưu hóa tuyển dụng nhân sự đến việc hỗ trợ HĐQT ra quyết định chiến lược dựa trên dữ liệu.

Thứ tư, bên cạnh những thành tựu rực rỡ, việc quy mô nhân sự vượt ngưỡng 70.000 người cũng bộc lộ những điểm nghẽn về độ trễ vận hành, hiện tượng cát cứ dữ liệu (silo) giữa các đơn vị thành viên, áp lực cạnh tranh nhân tài công nghệ cao và thách thức chuyển giao thế hệ lãnh đạo. Những nhóm giải pháp về tinh gọn bộ máy, xây dựng nền tảng dữ liệu hợp nhất (DC5), hoàn thiện Khung Quản trị AI và hoạch định nhân sự kế thừa mà nhóm đề xuất sẽ là những gợi mở thiết thực giúp FPT tiếp tục giữ vững vị thế tiên phong.

Nghiên cứu về FPT không chỉ giúp sinh viên củng cố sâu sắc hệ thống kiến thức môn Khoa học Quản lý và Quản trị, mà còn mang lại những bài học thực tiễn vô giá về tinh thần đổi mới sáng tạo, năng lực tổ chức bộ máy và tầm nhìn chiến lược cho các nhà quản lý tương lai.
